\section{Protocol elements and model}
\label{sec:model}

\subsection{JAM auditing and disputes}

The current Gray Paper states that JAM's auditing and judging system is theoretically equivalent to ELVES \citep{wood2026,burdges2024}. Each validator initially audits a random deterministic selection of up to ten active cores. If an announced auditor fails to produce a positive judgment, later tranches recruit additional auditors. The Gray Paper sets the audit bias factor to
\[
F=2,
\]
meaning that one unmatched announcement induces an expected two additional auditors in the next tranche. A negative judgment instead causes broad re-auditing and opens the dispute path.

The dispute rules are central to our first result. Let
\[
K=\floor{\frac{2N}{3}}+1
\]
be the number of judgments carried by a verdict and
\[
W=\floor{\frac{N}{3}}
\]
the number of positive judgments required for a wonky verdict. Under the current Gray Paper rule, a verdict is:
\begin{itemize}
\item \textbf{good} if all \(K\) included judgments are positive;
\item \textbf{bad} if all \(K\) included judgments are negative;
\item \textbf{wonky} if exactly \(W\) of the \(K\) judgments are positive.
\end{itemize}

Bad and wonky verdicts remove the availability assignment. The Gray Paper's ordinary offense paths are tied to good or bad reports: guarantors of a bad report can be included as culprits, and judges contradicting an established good/bad verdict can be included as faults. A wonky report is recorded in the wonky set, but the current specification does not place its guarantors or judges in the punish set through those culprit/fault rules \citep{wood2026}.

This paper therefore uses the precise phrase \emph{protocol-defined offender attribution path} rather than equating a wonky verdict with zero economic consequences. A staking subsystem could in principle react to the wonky record; the Gray Paper leaves validator rewards and much of the higher-level economics outside the JAM chain itself.

\subsection{ELVES branching quantities}

ELVES models the audit-tranche process with a pessimistic adaptive-crash adversary. Let \(s_0\) be the expected initial audit committee and let \(F\) be the expected number of auditors recruited per unmatched prior auditor. If a fraction \(\gamma\) of the validator population is adversarial, the mean number of useful honest offspring in the branching approximation is
\[
\lambda=(1-\gamma)F.
\]
If \(q\) is the extinction probability of the corresponding Poisson Galton--Watson process, then
\[
q=\exp[-\lambda(1-q)],
\]
and the ELVES soundness analysis yields the pessimistic acceptance bound
\[
\boxed{
P_{\rm accept}\le q^{s_0/F}.
}
\]
Useful escalation requires \(\lambda>1\).

The same architecture has a no-show/scalability side. If \(u\) is the fraction represented in the no-show/fault-amplification term, the reproduction number is \(uF\), so the static expected-final-committee bound has the form
\[
\boxed{
 s_f\le\frac{s_0}{1-uF},
 \qquad uF<1.
}
\]
ELVES reports a continuous-model optimum near \(F=2.0794\) under its stated assumptions; the Gray Paper uses \(F=2\).

\subsection{Correlated honest fault domain}

For one crafted invalid report and one fault class \(x\), partition the verdict population into:
\begin{align*}
A &: \text{adversarial validators},\\
B_x &: \text{otherwise-honest validators that judge the invalid report as valid},\\
C_x &: \text{otherwise-honest validators that independently judge it as invalid}.
\end{align*}
Thus
\[
N=A+B_x+C_x.
\]

The central distinction is:
\[
\boxed{
\text{honest validator count}
\neq
\text{fault-specific independently corrective validator count}.
}
\]

When convenient we write
\[
\gamma=\frac{A}{N},
\qquad
f_x=\frac{B_x}{N-A},
\]
so
\[
\frac{C_x}{N}=(1-\gamma)(1-f_x).
\]
The parameter \(f_x\) should not be interpreted as a universal client market share. It is the exposure of honest validators to fault \(x\). A client implementation is one possible fault domain, but shared libraries, compilers, virtual machines, cryptographic code, erasure-coding implementations, or copied specification logic can couple nominally distinct clients.

\subsection{Claim-status map}

\begin{table}[htbp]
\centering
\small
\caption{Status of the main claims. Extension means derived after changing an ELVES/JAM assumption; it is not attributed to the original authors.}
\label{tab:claims}
\begin{tabularx}{\textwidth}{@{}p{0.05\textwidth}X p{0.23\textwidth}@{}}
\toprule
ID & Claim & Status \\
\midrule
A1 & A bad verdict remains constructible against adversarial positive voting iff \(C_x\ge K\). & Exact consequence of Gray Paper verdict rule + declared fault model \\
A2 & A good verdict becomes constructible if \(A+B_x\ge K\). & Exact consequence of Gray Paper verdict rule + declared fault model \\
A3 & For large \(N\), robust bad-verdict attribution requires \(\gamma+(1-\gamma)f_x<1/3\). & Asymptotic form of A1 \\
B1 & Correlated honest failure changes the ELVES corrective mean to \(\lambda_x=F C_x/N\). & ELVES extension \\
B2 & Supercritical correction requires \(f_x<1-1/[(1-\gamma)F]\). & Exact consequence of B1 \\
B3 & Increasing \(F\) raises fault-domain tolerance while worsening the static no-show committee bound. & ELVES extension / comparative result \\
C1 & Holding \(C_x\) fixed, adding validators outside the independently correct set strictly lowers \(\lambda_x=FC_x/N\). & Exact sampling identity \\
C2 & In the declared baseline, the pessimistic acceptance bound worsens under such additions both for fixed \(s_0=30\) and fixed 341 cores. & Computational specialization \\
D1 & Slashable collateral need not be attacker-attributable collateral under an honest shared bug. & Economic mapping observation; not a theorem of JAM \\
D2 & Guarantor diversity, effort proofs, announcement monitoring and load feedback are design questions, not established fixes. & Hypotheses / future work \\
\bottomrule
\end{tabularx}
\end{table}
